\documentclass{article}

\usepackage[T1]{fontenc}

\usepackage{booktabs,graphicx,siunitx}

\begin{document}

\section{Composite reliability (CR)}

\begin{tabular}{lllll}
\toprule
Construct & CR & AVE & $\omega$ & $\alpha$ \\
\midrule
esg & .86 & .60 & .86 & .86 \\
norm & .88 & .65 & .88 & .88 \\
intent & .84 & .65 & .84 & .84 \\
attitude & .88 & .65 & .88 & .88 \\
service\_quality & .86 & .62 & .86 & .86 \\
\bottomrule
\end{tabular}

\includegraphics{figures/01_composite-reliability/figure_reliability.png}

CR (Composite Reliability) measures how reliably a set of items captures its construct (0--1); $\ge$ .70 is the conventional threshold (Nunnally, 1978; Bagozzi \& Yi, 1988). CR is computed from the CFA loadings via semTools::compRelSEM() - for a congeneric (unidimensional) factor CR equals McDonald's $\omega$, so the two columns will always match. AVE is shown alongside for quick convergent-validity reference ($\ge$ .50; Fornell \& Larcker, 1981) - for the full discriminant-validity assessment (Fornell--Larcker + HTMT matrices) run the AVE card. $\alpha$ (Cronbach's) is a secondary/legacy coefficient - it assumes tau-equivalence (equal loadings) and is a lower bound when loadings differ (McNeish, 2018). All cutoffs are labelled guidelines, not pass/fail gates. Applies to reflective constructs only.

Composite reliability for esg was satisfactory (CR = .86 $\ge$ .70).

Statistical basis: CR (composite reliability) $\ge$ .70 acceptable (Nunnally, J. C. (1978). Psychometric Theory (2nd ed.). McGraw-Hill.); CR $\ge$ .70 threshold, independent corroboration (Bagozzi, R. P., Yi, Y. (1988). "On the evaluation of structural equation models." Journal of the Academy of Marketing Science, 16(1), 74-94.); AVE shown alongside for convergent-validity reference ($\ge$ .50) (Fornell, C., Larcker, D. F. (1981). "Evaluating structural equation models with unobservable variables and measurement error." Journal of Marketing Research, 18(1), 39-50.); $\alpha$ retained as a secondary/legacy coefficient (lower bound when loadings differ) (McNeish, D. (2018). "Thanks coefficient alpha, we'll take it from here." Psychological Methods, 23(3), 412-433.).

\end{document}